\documentclass[11pt]{article}

\usepackage[hidelinks]{hyperref}
\usepackage{xcolor}
\usepackage{soul}
\usepackage{amsmath}
\usepackage{amssymb}
\usepackage{mathtools}
\usepackage{gensymb}
\usepackage{graphicx}
\usepackage{caption}
\usepackage{wrapfig}
\usepackage{tabularx}
\usepackage{booktabs}
\usepackage{multirow}

\title{Executive Summary for the U.S Defense Report}
\author{Overview of  U.S Defense Infrastructure Investment in the Philippines: Funding, Geographic Access, and Infrastructure Delivery}
\date{Research period FY2016 - FY2023}

\begin{document}

\section{Purpose }
The expansion of the Enhanced Defense Cooperation Agreement (EDCA) from five to nine Philippine locations has increased the geographic distribution of U.S military access. However publicly available site-level funding information remains incomplete.

This research identifies approximately \$86.17 million in reported allocation across five original EDCA locations, with approximately 90.5\% concentrated at Basa Air Base and Fort Magsaysay. Because the available records do not established actual expenditures, completed construction, or operational availability across all nine locations, the full extended infrastructure investment and delivery cannot be independently.

\end{document}
